\section{Bivariate Analysis: Is X Related to Y?}

Bivariate analysis is used to analyze two variables simultaneously. In
this study, the two variables of interest are:

\[
X = \text{Use of Gadgets as a Learning Medium}
\]
\[
Y = \text{Students' Learning Motivation and Social Interaction}
\]

The study tests the null hypothesis ($H_0$) that there is no relationship
between X and Y against the alternative hypothesis ($H_1$) that such a
relationship exists.

One way to examine their relationship is by calculating the correlation
coefficient.

The correlation coefficient ranges from $-1$ to $+1$. A positive value
indicates that the variables tend to move in the same direction, while a
negative value indicates an opposite direction. The closer the value is
to $+1$ or $-1$, the stronger the observed linear relationship.

\subsection{Pearson Correlation}

The Pearson correlation coefficient obtained in this study was:
\[
r = \resR
\]
This indicates a strong positive relationship between variables X and Y.
In other words, respondents with higher gadget-use scores tended to have
higher scores for learning motivation and social interaction. The
correlation matrix is visualized in Figure~\ref{fig:heatmap}.

\begin{figure}[htbp]
  \centering
  \includegraphics[width=0.5\linewidth]{figures/correlation_heatmap.pdf}
  \caption{Correlation matrix for gadget use (X) and learning motivation
  and social interaction (Y).}
  \label{fig:heatmap}
\end{figure}

The significance value was:
\[
\text{Sig.} = \resSig < 0.05
\]
This indicates that the observed relationship was statistically
significant within the tested sample and model.

However, statistical significance should not automatically be interpreted
as proof that gadget use causes changes in learning motivation or social
interaction. The research design is correlational, meaning that the
analysis identifies a relationship between variables rather than
establishing a definitive cause-and-effect relationship.

\subsection{R-Square ($R^2$)}

The analysis also produced an $R^2$ value of \resRSquared{}.

\subsection{Calculation 2 --- Converting $R^2$ into a Percentage}

Formula:
\[
R^2 \times 100\%
\]
\[
\resRSquared \times 100\% = \resRSquaredPct\%
\]
This means that approximately \resRSquaredPct\% of the variation in Y can
be explained by X within the research model.

The remaining variation can be calculated as:
\[
(1 - \resRSquared) \times 100\% = \resUnexplainedPct\%
\]
Therefore, approximately \resUnexplainedPct\% of the variation in Y is not
explained by X within the model and may be associated with other factors
that were not included in the analysis.

This result shows that X has a substantial statistical relationship with
Y in the analyzed model. However, X is not the only factor associated
with variation in Y.

\subsection{t-Test}

The analysis also produced $t = \resT$ with a significance value of
\resSig{}.

Because:
\[
\resSig < 0.05
\]
the result indicates that variable X has a statistically significant
relationship with variable Y within the tested regression model.

Therefore:
\[
\text{Sig. } \resSig < 0.05 \rightarrow H_0 \text{ is rejected and } H_1 \text{ is accepted.}
\]

\subsection{F-Test}

The F-test produced $F = \resF$ with a significance value of \resSig{}.

Because:
\[
\resSig < 0.05
\]
the result indicates that the regression model as a whole is
statistically significant.

Therefore:
\[
\text{Sig. } \resSig < 0.05 \rightarrow \text{the research model is statistically significant.}
\]

\subsection{Bivariate Analysis Conclusion}

Based on the results, the alternative hypothesis ($H_1$) is accepted. The
analysis indicates a positive and statistically significant relationship
between variables X and Y in the sample and model studied. This finding
is supported by the Pearson correlation coefficient of \resR{} with a
significance value of \resSig{}.

The $R^2$ value of \resRSquaredPct\% indicates that \resRSquaredPct\% of
the variation in Y
can be explained by X within the research model, while the remaining
\resUnexplainedPct\% is not explained by X within the model.

Nevertheless, the statistical results need to be interpreted carefully. A
statistically significant relationship does not automatically mean that X
is the sole cause of changes in Y. Other factors may also be related to
variation in Y.
